% ============================================================
% 3.4.2 RISK EVALUATION
% ============================================================

\hypertarget{p3-4-2}{}

\subsection{Risk Evaluation}

\subsubsection{General}

The risk evaluation of the \textbf{BN-100 LUX} Infant Radiant Warmer is
addressed through the manufacturer's risk management documentation.

The evaluation considers the risks associated with the identified hazards
and hazardous situations and their relationship to the intended use and
safety-related functions of the device.

The available Risk Management Report provides the documentary basis for the
risk evaluation.

\subsubsection{Risk Evaluation Basis}

The manufacturer's risk management documentation is identified as:

\begin{itemize}

    \item \textbf{Risk Management Report (BN-100/BN-100A)};

    \item Document No.: \textbf{YF-7.3-24-02 (w)};

    \item Version: \textbf{A/0}.

\end{itemize}

The document title is retained as provided in the manufacturer
documentation. The device covered by the present dossier remains the
\textbf{BN-100 LUX}.

\subsubsection{Evaluation of Identified Risks}

The identified risks are evaluated according to their potential effect on
the safety and performance of the device.

The risk evaluation takes into consideration the nature of the hazard, the
reasonably foreseeable conditions that may lead to a hazardous situation and
the potential harm to the patient, operator or other persons.

The manufacturer's risk management process provides the basis for
determining whether the identified risks require risk-control measures.

\subsubsection{Electrical Risk Evaluation}

Electrical hazards are evaluated in relation to the electrical architecture
and components of the device.

Relevant components include the power supply, input socket, fuse, relay,
switching power supply, filter, internal wiring and printed circuit board.

The evaluation considers the possible consequences of electrical component
failure or abnormal electrical conditions and their potential impact on
patient and operator safety.

\subsubsection{Thermal Risk Evaluation}

Thermal risks are particularly relevant to the \textbf{BN-100 LUX} because
the principal function of the device is to provide controlled radiant heat.

The evaluation therefore considers conditions such as:

\begin{itemize}

    \item excessive heating;

    \item insufficient or uncontrolled heating;

    \item temperature-control failure;

    \item temperature-monitoring failure;

    \item deviation from the specified temperature condition.

\end{itemize}

These risks are considered together with the documented temperature-control
and alarm functions of the device.

\subsubsection{Temperature Sensor Risk Evaluation}

The skin temperature sensor is a safety-related component of the
temperature-monitoring system.

The manufacturer documentation identifies sensor-related failure conditions,
including short circuit and open circuit.

These conditions are evaluated in relation to their potential effect on
temperature monitoring and patient safety.

The corresponding alarm functions provide a documented means of detecting
specified sensor failures.

\subsubsection{Alarm System Risk Evaluation}

The manufacturer identifies several alarm conditions associated with
potentially hazardous operating conditions.

These include over-temperature, temperature deviation, sensor failure,
power failure and manual-mode operation.

The evaluation of these conditions is related to the ability of the alarm
system to provide an indication of the corresponding abnormal condition.

\subsubsection{Mechanical Risk Evaluation}

Mechanical risks are evaluated in relation to the structure and intended use
of the patient support system and associated components.

The manufacturer's documented load limits for the mattress support and tray
provide defined mechanical operating limits.

The crib inclination and adjustment mechanism are also considered in
relation to the intended positioning function of the patient support.

\subsubsection{Patient-Contact Risk Evaluation}

The risk evaluation includes consideration of components that are in direct
contact with the patient.

The manufacturer identifies the cotton baby mattress and stainless-steel
skin temperature sensor probe as patient-contact components.

The corresponding biological considerations are addressed within the
manufacturer's risk management documentation.

\subsubsection{Manufacturing Risk Evaluation}

The manufacturer's risk management documentation considers risks that may
arise from manufacturing and assembly processes.

Such risks are evaluated according to their potential effect on the safety
and performance of the finished device.

The corresponding process controls and manufacturing controls form part of
the overall risk-control framework.

\subsubsection{Risk Acceptability}

The acceptability of identified risks is determined within the
manufacturer's documented risk management process.

This section does not establish independent numerical acceptance criteria or
introduce a new risk matrix.

Where detailed severity, probability, risk levels or acceptance decisions
are required, the original manufacturer's Risk Management Report shall be
used as the controlling documentary reference.

\subsubsection{Overall Evaluation Approach}

The risk evaluation considers the relationship between the identified
hazards, hazardous situations, potential harms and the corresponding
risk-control measures.

The evaluation is performed in the context of the intended use of the
\textbf{BN-100 LUX} and its documented safety-related functions.

The available manufacturer documentation therefore provides the primary
basis for the evaluation of the identified risks.

\subsubsection{Documentation Limitation}

The present section summarizes the risk evaluation documented by the
manufacturer.

No independent severity, probability or risk-ranking values are introduced
where they are not explicitly available in the supporting manufacturer
documentation.

The original risk-analysis and risk-evaluation records should be retained
as part of the technical documentation supporting this section.

\subsubsection{Documentary Basis}

This section is based primarily on:

\begin{itemize}

    \item Risk Management Report (BN-100/BN-100A), Document No.
    YF-7.3-24-02 (w), Version A/0;

    \item BN-100 / BN-100 LUX User Manual.

\end{itemize}
